\par\Needspace{6\baselineskip}\medskip
\begin{center}\hypertarget{native-185-article138}{}\pdfbookmark[2]{138.}{article-138}138.\end{center}
Die \emph{dritte} Aufgabe besteht nun darin, den ersten und den letzten geozentrischen Ort des
Himmelskörpers, d.\,h. die Punkte $B,B''$, durch einen Großkreis zu verbinden und dessen
Schnittpunkt mit dem Großkreis $A'B'$ zu bestimmen. Es sei $B^*$ dieser Schnittpunkt, $\delta'-\sigma$ sein
Abstand vom Punkt $A'$, $\alpha^*$ seine Länge und $\beta^*$ seine Breite. Da $B,B^*,B''$ auf demselben
Großkreis liegen, gilt die bekannte Gleichung
\[
\begin{aligned}
0={}&\operatorname{tang}\beta\sin(\alpha''-\alpha^*)-\operatorname{tang}\beta^*\sin(\alpha''-\alpha)\\
&+\operatorname{tang}\beta''\sin(\alpha^*-\alpha),
\end{aligned}
\]
Setzt man darin $\operatorname{tang}\gamma'\sin(\alpha^*-l')$ für $\operatorname{tang}\beta^*$ ein, so nimmt sie die folgende Form an
\[
\begin{aligned}
0={}&\cos(\alpha^*-l')\{\operatorname{tang}\beta\sin(\alpha''-l')-\operatorname{tang}\beta''\sin(\alpha-l')\}\\
&-\sin(\alpha^*-l')\{\operatorname{tang}\beta\cos(\alpha''-l')\\
&\qquad+\operatorname{tang}\gamma'\sin(\alpha''-\alpha)-\operatorname{tang}\beta''\cos(\alpha-l')\}.
\end{aligned}
\]
Da nun $\operatorname{tang}(\alpha^*-l')=\cos\gamma'\operatorname{tang}(\delta'-\sigma)$ gilt, erhält man
\[
\begin{gathered}
\operatorname{tang}(\delta'-\sigma)=\\[.4ex]
\frac{\operatorname{tang}\beta\sin(\alpha''-l')-\operatorname{tang}\beta''\sin(\alpha-l')}{\begin{gathered}\cos\gamma'(\operatorname{tang}\beta\cos(\alpha''-l')-\operatorname{tang}\beta''\cos(\alpha-l'))\\+\sin\gamma'\sin(\alpha''-\alpha)\end{gathered}}\mathpunct{\cdot}
\end{gathered}
\]

Daraus leiten wir die folgenden, für die numerische Rechnung besser geeigneten Formeln ab. Man
setze
\[
\begin{aligned}
7)\quad&\operatorname{tang}\beta\sin(\alpha''-l')-\operatorname{tang}\beta''\sin(\alpha-l')=S\\[.5ex]
8)\quad&\operatorname{tang}\beta\cos(\alpha''-l')-\operatorname{tang}\beta''\cos(\alpha-l')=T\sin t\\[.5ex]
9)\quad&\sin(\alpha''-\alpha)=T\cos t
\end{aligned}
\]
(Artikel 14, II). Dann ergibt sich
\[
10)\quad\operatorname{tang}(\delta'-\sigma)=\frac{S}{T\sin(t+\gamma')}\mathpunct{\cdot}
\]
Die Bestimmung des Bogens $\delta'-\sigma$ aus seinem Tangens ist mehrdeutig, weil sich die Großkreise
$A'B',BB''$ in \emph{zwei} Punkten schneiden. Für $B^*$ wählen wir stets den \hypertarget{native-186}{}Schnittpunkt, der
dem Punkt $B'$ am nächsten liegt, sodass $\sigma$ stets zwischen den Grenzen $-90^0$ und $+90^0$
liegt. Damit ist die Mehrdeutigkeit beseitigt. Der Wert des Bogens $\sigma$ (der von der
\emph{Krümmung} der geozentrischen Bewegung abhängt) ist dann meist recht klein. Im Allgemeinen
ist er von zweiter Ordnung, wenn man die Zeitabstände als Größen erster Ordnung betrachtet.

Welche Änderungen an der Rechnung erforderlich sind, wenn man als Grundebene anstelle der
Ekliptik den Äquator wählt, ergibt sich unmittelbar aus der Bemerkung im vorigen Artikel.

Offenbar bliebe die Lage des Punkts $B^*$ unbestimmt, wenn die Kreise $BB'',A'B'$ vollständig
zusammenfielen. Diesen Fall, in dem die vier Punkte $A',B,B',B''$ auf demselben Großkreis lägen,
schließen wir von unserer Untersuchung aus. Bei der Wahl der Beobachtungen empfiehlt es sich
auch, den Fall zu vermeiden, dass die vier Punkte nahe bei einem einzigen Großkreis liegen. Dann
würden schon geringste Beobachtungsfehler die Lage des Punkts $B^*$, auf die es bei den
folgenden Rechenschritten wesentlich ankommt, zu stark beeinflussen. Sie ließe sich daher nicht
mit der erforderlichen Genauigkeit bestimmen. --- Ebenso bleibt der Punkt $B^*$ offenbar
unbestimmt, wenn die Punkte $B,B''$ zusammenfallen\footnote{Oder wenn sie einander
gegenüberliegen. Diesen Fall erörtern wir nicht, da unser Verfahren nicht auf Beobachtungen
ausgedehnt werden soll, die einen so langen Zeitraum umfassen.}, denn dann wäre die Lage des
Kreises $BB''$ selbst unbestimmt. Deshalb schließen wir auch diesen Fall aus. Aus ähnlichen
Gründen wie zuvor sind auch Beobachtungen zu vermeiden, bei denen der erste und der letzte
geozentrische Ort auf der Kugel nahe beieinanderliegen.
